\begin{table}[htbp]
\centering
\caption{Placebo IV: Sub-Threshold Always-Losers as Instrument}
\label{tab:iv_placebo}
\begin{adjustbox}{max width=\textwidth}
\begin{threeparttable}
\small
\begin{tabular}{lcc}
\toprule
 & First Stage & 2SLS \\
 & FL presence & Log Price \\
\midrule
Sub-threshold supply (LOO) & 0.000040*** & \\
 & (0.000007) & \\
F-statistic & 33.0 & \\
FL presence (instrumented) & & 5.2445*** \\
 & & (1.3200) \\
\bottomrule
\end{tabular}
\begin{tablenotes}
\small
\item \textit{Notes:} Placebo instrument uses supply of sub-threshold
always-losers (below the IQR cutoff). The sub-threshold first stage
($F = 33$) exceeds the \citet{stock2005testing} 10\% maximal IV size
critical value of 16.38 for a single instrument. The implausibly large
2SLS coefficient (5.24) is therefore not attributable to weak-instrument
bias. Rather, it reflects the absence of a genuine cover-bidding
mechanism: sub-threshold always-losers do not drive procurement price
anomalies, confirming that the main instrument's price effect operates
through the cover-bidding channel rather than through correlated demand
shocks. SE clustered at item level.
\end{tablenotes}
\end{threeparttable}
\end{adjustbox}
\end{table}
